\chapter{1977年安徽卷（文）}

\section{解答题}

\begin{problem}
计算：

\begin{enumerate}
    \item \(5\frac{1}{2}\times\left(-\frac{6}{11}\right)+0.25-(-2)^3\div\left(-2\frac{2}{3}\right)^2\)；
    \item \(\sqrt{14\times24\times39\times91}+4\sqrt{27}\times(-\sqrt{3})\)；
    \item \(\left(-2\frac{3}{5}\right)^0+4^{-2}\times\left(2\frac{1}{4}\right)^{-\frac{1}{2}}-(0.01)^{0.5}\)；
    \item \(\frac{2\lg2+\lg3}{2+\lg0.36+\frac{2}{3}\lg8}\)；
    \item \(\frac{\tan\left(-120^\circ\right)\cos\left(-240^\circ\right)\cos480^\circ}{\tan\left(-60^\circ\right)\sin\left(-1050^\circ\right)}\).
\end{enumerate}
\end{problem}

\begin{solution}
\begin{enumerate}
    \item 先把带分数和小数化为分数，并注意减去负数所得的符号，原式 \(=\frac{11}{2}\times\left(-\frac{6}{11}\right)+\frac{1}{4}-\frac{-8}{\left(-\frac{8}{3}\right)^2}=-3+\frac{1}{4}+\frac{9}{8}=-\frac{13}{8}=-1\frac{5}{8}\).

    \item 将各因数分解并配对，得 \(14=(2\times7)\)、\(24=(3\times8)\)、\(39=(3\times13)\)、\(91=(7\times13)\)，所以
    \(\sqrt{14\times24\times39\times91}=\sqrt{(2\times7)(3\times8)(3\times13)(7\times13)}=\sqrt{(4\times3\times7\times13)^2}=4\times3\times7\times13=1092\). 又 \(4\sqrt{27}\times(-\sqrt3)=4\times3\sqrt3\times(-\sqrt3)=-36\)，故原式 \(=1092-36=1056\).

    \item 因为 \(-2\frac{3}{5}\neq0\)，所以 \(\left(-2\frac{3}{5}\right)^0=1\)，且 \(4^{-2}=\frac{1}{16}\)，\(\left(2\frac{1}{4}\right)^{-\frac12}=\left(\frac94\right)^{-\frac12}=\frac{1}{\sqrt{9/4}}=\frac23\)，\((0.01)^{0.5}=\sqrt{0.01}=\frac{1}{10}\). 因此原式 \(=1+\frac{1}{16}\times\frac23-\frac{1}{10}=1+\frac{1}{24}-\frac{1}{10}=\frac{113}{120}\).

    \item 由 \(0.36=\frac{36}{100}\) 及 \(8=2^3\)，得 \(\lg0.36=\lg36-2=2\lg2+2\lg3-2\)，\(\frac23\lg8=2\lg2\). 因而分母 \(=2+(2\lg2+2\lg3-2)+2\lg2=2(2\lg2+\lg3)\). 又 \(2\lg2+\lg3=\lg(2^2\cdot3)=\lg12>0\)，分子非零，故原式 \(=\frac{1}{2}\).

    \item 利用诱导公式，\(\tan(-120^\circ)=\tan60^\circ\)，\(\cos(-240^\circ)=-\cos60^\circ\)，\(\cos480^\circ=-\cos60^\circ\)，\(\tan(-60^\circ)=-\tan60^\circ\)，\(\sin(-1050^\circ)=\sin30^\circ\). 所以原式 \(=\frac{\tan60^\circ\cos^2 60^\circ}{-\tan60^\circ\sin30^\circ}=-\frac{\cos^2 60^\circ}{\sin30^\circ}=-\frac{(1/2)^2}{1/2}=-\frac12\).
\end{enumerate}
\end{solution}

\begin{problem}
化简

\begin{enumerate}
    \item \(\left(3\sqrt[3]{a^2}\cdot\sqrt b\right)\left(-6\sqrt a\cdot\sqrt[3]b\right)\div\left(-2\sqrt[6]a\cdot\sqrt[6]{b^5}\right)\)；
    \item \(\sqrt{(\lg x)^2-2\lg x+1}\).
\end{enumerate}
\end{problem}

\begin{solution}
\begin{enumerate}
    \item 原式有意义且分母不为零时，根式要求 \(a\geqslant0\)、\(b\geqslant0\)，分母不为零进一步要求 \(a>0\)、\(b>0\). 在此条件下，利用根式的幂指数表示，原式 \(=3a^{2/3}b^{1/2}\cdot(-6a^{1/2})b^{1/3}\div\left[(-2a^{1/6})b^{5/6}\right]\). 合并系数与同底幂，并计算指数，得原式 \(=9a^{2/3+1/2-1/6}b^{1/2+1/3-5/6}=9a^{4/6+3/6-1/6}b^{3/6+2/6-5/6}=9a^1b^0=9a\).

    \item 先利用对数的定义域知 \(x>0\)，再配方，原式 \(=\sqrt{(\lg x-1)^2}=|\lg x-1|\). 若写成分段形式，则 \(\lg x-1\geqslant0\) 等价于 \(x\geqslant10\)，所以
    \[
    \sqrt{(\lg x-1)^2}=\begin{cases}
    1-\lg x,&0<x\leqslant10,\\
    \lg x-1,&x>10.
    \end{cases}
    \]
\end{enumerate}
\end{solution}

\begin{problem}
分解因式：\(a^2-2ab+b^2-6a+6b+5\).
\end{problem}

\begin{solution}
先把同类项按 \(a-b\) 组合，有 \(a^2-2ab+b^2=(a-b)^2\)、\(-6a+6b=-6(a-b)\). 令 \(u=a-b\)，则原式 \(=u^2-6u+5=(u-1)(u-5)\). 代回 \(u=a-b\)，得到原式 \(=(a-b-1)(a-b-5)\).
\end{solution}

\begin{problem}
求下列函数的定义域，并在数轴上把它们表示出来：

\begin{enumerate}
    \item \(y=\frac{\sqrt{3+2x}}{3x-1}\)；
    \item \(y=\arcsin\frac{x}{3}\).
\end{enumerate}
\end{problem}

\begin{solution}
\begin{enumerate}
    \item 根式要求 \(3+2x\geqslant0\)，即 \(x\geqslant-\frac32\)；分母要求 \(3x-1\neq0\)，即 \(x\neq\frac13\). 两个条件取交集，定义域为 \(\left[-\frac32,\frac13\right)\cup\left(\frac13,+\infty\right)\). 数轴上在 \(-\frac32\) 处取闭点，在 \(\frac13\) 处取空心点，并标出从 \(-\frac32\) 到 \(\frac13\) 以及从 \(\frac13\) 向右的两段.

    \item 反正弦函数的自变量必须属于 \([-1,1]\)，故 \(-1\leqslant\frac{x}{3}\leqslant1\). 由于 \(3>0\)，不等式同向乘以 \(3\) 得 \(-3\leqslant x\leqslant3\). 因此定义域为 \([-3,3]\)，数轴上在 \(-3\)、\(3\) 两端均取闭点，并标出两点之间的线段.
\end{enumerate}
\end{solution}

\begin{problem}
求证：\(\frac{\sin\alpha+\sin\beta}{\sin(\alpha+\beta)}=\frac{\cos\frac{\alpha-\beta}{2}}{\cos\frac{\alpha+\beta}{2}}\).
\end{problem}

\begin{solution}
在等式两边有意义的条件下，左边的分母要求 \(\sin(\alpha+\beta)\neq0\). 由 \(\sin(\alpha+\beta)=2\sin\frac{\alpha+\beta}{2}\cos\frac{\alpha+\beta}{2}\)，可知 \(\sin\frac{\alpha+\beta}{2}\neq0\) 且 \(\cos\frac{\alpha+\beta}{2}\neq0\)，因此下面的约分合法. 由和差化积公式，\(\sin\alpha+\sin\beta=2\sin\frac{\alpha+\beta}{2}\cos\frac{\alpha-\beta}{2}\)，且 \(\sin(\alpha+\beta)=2\sin\frac{\alpha+\beta}{2}\cos\frac{\alpha+\beta}{2}\)，所以
    \[
    \frac{\sin\alpha+\sin\beta}{\sin(\alpha+\beta)}=\frac{2\sin\frac{\alpha+\beta}{2}\cos\frac{\alpha-\beta}{2}}{2\sin\frac{\alpha+\beta}{2}\cos\frac{\alpha+\beta}{2}}=\frac{\cos\frac{\alpha-\beta}{2}}{\cos\frac{\alpha+\beta}{2}}
    \]
    这正是右端，故原等式成立.
\end{solution}

\begin{problem}
解方程 \(\sin x+\cos x-1=0\).
\end{problem}

\begin{solution}
由原方程得 \(\sin x=1-\cos x\). 再用半角公式 \(\sin x=2\sin\frac{x}{2}\cos\frac{x}{2}\) 和 \(1-\cos x=2\sin^2\frac{x}{2}\)，得到 \(2\sin\frac{x}{2}\cos\frac{x}{2}=2\sin^2\frac{x}{2}\)，即 \(\sin\frac{x}{2}\left(\cos\frac{x}{2}-\sin\frac{x}{2}\right)=0\).

    因此分两种情况：当 \(\sin\frac{x}{2}=0\) 时，\(\frac{x}{2}=k\pi\)，得 \(x=2k\pi\)；当 \(\cos\frac{x}{2}-\sin\frac{x}{2}=0\) 时，\(\cos\frac{x}{2}\neq0\)，从而 \(\tan\frac{x}{2}=1\)，得 \(\frac{x}{2}=k\pi+\frac\pi4\)，即 \(x=2k\pi+\frac\pi2\)，其中 \(k\in\Z\). 回代检验：当 \(x=2k\pi\) 时，\(\sin x+\cos x-1=0+1-1=0\)；当 \(x=2k\pi+\frac\pi2\) 时，\(\sin x+\cos x-1=1+0-1=0\). 因此方程的解为 \(x=2k\pi\) 或 \(x=2k\pi+\frac\pi2\)，\(k\in\Z\).
\end{solution}

\begin{problem}
求下列等式里的实数 \(x\) 和 \(y\) 的值：\((3x\i+2x)+(y\i-4)=(2y+5\i)-(3+x\i)\).
\end{problem}

\begin{solution}
将等式两边分别合并实部和虚部，左边为 \((2x-4)+(3x+y)\i\)，右边为 \((2y-3)+(5-x)\i\). 复数相等当且仅当实部、虚部分别相等，因此 \(2x-4=2y-3\)，\(3x+y=5-x\)，即
\[
\begin{cases}
2x-2y=1,\\
4x+y=5.
\end{cases}
\]
将第一式与第二式的 2 倍相加，得 \(10x=11\)，所以 \(x=\frac{11}{10}\). 代入第二式，得 \(y=5-4\times\frac{11}{10}=\frac35\). 检验可得两边实部均为 \(-\frac95\)，虚部均为 \(\frac{39}{10}\)，故 \(x=\frac{11}{10}\)，\(y=\frac35\).
\end{solution}

\begin{problem}
从山顶 \(D\) 测得地面上同一方向的两点 \(A\) 和 \(B\) 的俯角分别是 \(30^\circ\) 和 \(45^\circ\)，已知 \(AB=40\text{ 米}\)，求山高. （精确到 \(0.1\)）
\end{problem}

\begin{solution}
\solutionfigure{\bitmapfigure[max width=6cm]{img/1977/anhui_liberal/q8_fig1.png}}

设山脚为 \(C\)，则 \(DC\perp AC\)，且 \(A\)，\(B\)，\(C\) 在同一直线上，\(B\) 位于 \(A\)，\(C\) 之间. 俯角等于相应的仰角，所以 \(\angle DAC=30^\circ\)，\(\angle DBC=45^\circ\). 在直角三角形 \(DBC\) 中，\(\tan45^\circ=\frac{DC}{BC}=1\)，从而 \(BC=DC\)，于是 \(AC=AB+BC=40+DC\).

在直角三角形 \(DAC\) 中，\(\tan30^\circ=\frac{DC}{AC}=\frac{DC}{40+DC}\)，即 \(\frac{\sqrt3}{3}=\frac{DC}{40+DC}\). 交叉相乘得 \(\sqrt3(40+DC)=3DC\)，所以 \((3-\sqrt3)DC=40\sqrt3\)，从而 \(DC=\frac{40\sqrt3}{3-\sqrt3}=\frac{40}{\sqrt3-1}=20(\sqrt3+1)\approx54.641\text{ 米}\). 按题意精确到 \(0.1\)，山高为 \(54.6\text{ 米}\).
\end{solution}

\begin{problem}
红旗大队粮食产量逐年增加，1973年产量为 \(90\text{ 万斤}\)，连续三年平均每年比前一年增产 \(10\%\)，这个大队从1973年到1976年总共生产粮食多少万斤？（精确到 \(0.1\)）
\end{problem}

\begin{solution}
设 1973 年至 1976 年的产量分别为 \(a_0\)，\(a_1\)，\(a_2\)，\(a_3\)，单位均为万斤. 连续三年每年比前一年增产 \(10\%\)，所以相邻两年的产量比为 \(1+10\%=1.1\)，从而 \(a_0=90\)，\(a_1=90\times1.1=99\)，\(a_2=90\times1.1^2=108.9\)，\(a_3=90\times1.1^3=119.79\).

因此总产量 \(S=a_0+a_1+a_2+a_3=90(1+1.1+1.1^2+1.1^3)=90(1+1.1+1.21+1.331)=417.69\text{ 万斤}\). 按题意精确到 \(0.1\)，总产量为 \(417.7\text{ 万斤}\).
\end{solution}

\begin{problem}
求椭圆 \(25x^2+9y^2=100\) 的长轴和短轴的长，焦点坐标，并且画出它的图象.
\end{problem}

\begin{solution}
将方程两边除以 \(100\)，得 \(\frac{x^2}{4}+\frac{y^2}{100/9}=1\). 因为 \(\frac{100}{9}>4\)，所以椭圆的长轴在 \(y\) 轴上，长半轴为 \(a=\frac{10}{3}\)，短半轴为 \(b=2\). 因而长轴长为 \(2a=\frac{20}{3}\)，短轴长为 \(2b=4\).

半焦距满足 \(c^2=a^2-b^2\)，所以 \(c=\sqrt{\frac{100}{9}-4}=\frac83\). 因此两个焦点为 \(\left(0,\frac83\right)\) 和 \(\left(0,-\frac83\right)\). 作图时以原点为中心，取长轴端点 \(\left(0,\frac{10}{3}\right)\)、\(\left(0,-\frac{10}{3}\right)\) 和短轴端点 \((2,0)\)、\((-2,0)\)，作经过这些端点的竖直椭圆即可.

\solutionfigure{\bitmapfigure[max width=6cm]{img/1977/anhui_liberal/q10_fig1.png}}
\end{solution}
